\section{Protocolli DataLink}
(PDU dette frame).
\subsection{Controllo di canale}
\paragraph{Premessa} Tra i due estremi del mezzo trasmissivo i bit non cambiano l'ordine con cui sono stati trasmessi dalla sorgente. Ma a livello di pacchetto l'ordine può cambiare, per via dei nodi di rete.

\begin{paracol}{2}
	\paragraph{Protocolli di linea} è un canale sequenziale a banda costante 1-1 o 1-N in cui i frame frame arrivano nella stessa sequenza d'invio con ritardi di propagazione uguali.
	\switchcolumn
	\paragraph{Protocolli di trasporto} è un canale non sequenziale a banda variabile. Soffre di perdita di dati, duplicazione dei dati, ritardi variabili e arrivi fuori sequenza.
\end{paracol}

\paragraph{Scopo} Rendere affidabile e sicuro il servizio di collegamento offerto dallo strato 2 allo strato 3.

Funzioni svolte (non dallo stesso protocollo): strutturazione del flusso, controllo/gestione errori, controllo flusso, controllo sequenza, gestione protocollo accesso collegamento 1-N.

\paragraph{Problematiche di Sincronismo} In ricezione per riconoscere i bit bisogna determinare gli istanti di campionamento per ricostruire il sincronismo di cifra (problema riflesso dallo strato 1).

Il ricevitore estrae il segnale di sincronismo che può agganciarsi in due modi:
\begin{itemize}
	\item \textbf{Canale sempre pieno} mantenuto dal protocollo di linea mediante successivamente all'inizializzazione dell'aggancio sempre mantenuto.
	\item \textbf{Canale momentaneamente vuoto} con all'inizio di ogni tramissione l'inserimento di un preambolo di sincronizzazione.
\end{itemize}
\subparagraph{Sincronismo di cifra} garantisce corretta lettura singoli bit.
\subparagraph{Sincronismo di frame} distingue le PDU tramite:
\begin{itemize}
	\item Protocolli asincroni che permettono ai frame di iniziare e finire arbritariamente grazie info nel PCI per riconoscere inizio/fine.
	\item Protocolli sincroni che fanno iniziare e finire in istanti precisi i frame. Non necessitano di PCI.
\end{itemize}
\paragraph{Condifica di canale} come gli altri blocchi esegue una decodifica in ricezione.

\begin{adjustbox}{width=\textwidth}
	\begin{tikzcd}
		Sorgente \arrow[r] & \text{Codifica di Sorgente} \arrow[rr] & {} & {\text{[Codifica di Canale]}} \arrow[rr] & {}                                                                                                             & \text{Codifica di Linea} \arrow[rr] & {} & \text{Meezzo Fisico} \\
		&                                        &    &                                          & \text{Codifica Crittografica} \arrow[u, dashed] \arrow[rru, dashed, bend right] \arrow[llu, dashed, bend left] &                                     &    &
	\end{tikzcd}
\end{adjustbox}
\subsubsection{Controllo d'errore}
\paragraph{Codici a bloccho} codifica a blocchi di $k$ bit d'informazione. Vengono trasmessi $n=k+r$ bit dove $k$ sono i dati e $r$ i bit di rindondanza in funzione combinatoria di $k$.

I messaggi possibili sono $2^k$ e le parole di codice possibili sono $2^n>2^k$ di cui $2^n-2^k$ non sono lecite e usate per rilevare o rilevare e correggere errori.

Nota: Solo in binario è possibile fare la correzione.
\begin{paracol}{2}
	\subparagraph{Rivelazione d'errore} Una parola invalida indica la presenza d'un errore. Non si sanno i bit errati. Necessita della ritrasmissione dei dati.

	Migliore se faccio tanti errori.
	\switchcolumn
	\subparagraph{Correzione d'errore} Una parola invalida indica la presenza d'un errore e corregge il bit (trasparenza semantica).

	Migliore se faccio pochi errori.
\end{paracol}
